%% Esempio per lo stile supsi
\documentclass[twoside]{supsistudent}

%%%%%%%%%%%%%%%%%%%%%%%%%%%%%%%%%%%%%%%%%%%%%%%%%%%%%%%%%%%%
%%%   PACKAGES                                           %%%
%%%%%%%%%%%%%%%%%%%%%%%%%%%%%%%%%%%%%%%%%%%%%%%%%%%%%%%%%%%%
\usepackage{currfile}
\usepackage{hyperref}
\usepackage{float}
\usepackage{tabularx}
\usepackage{xcolor}
\usepackage{parskip}

\usepackage[linesnumbered, ruled, vlined]{algorithm2e}

\usepackage{enumitem}
\setlist[itemize]{noitemsep, nolistsep}



%%%%%%%%%%%%%%%%%%%%%%%%%%%%%%%%%%%%%%%%%%%%%%%%%%%%%%%%%%%%
%%%   COMMANDS                                           %%%
%%%%%%%%%%%%%%%%%%%%%%%%%%%%%%%%%%%%%%%%%%%%%%%%%%%%%%%%%%%%

%%% custom commands %%%%%%%%%%%%%%%%%%%%%%%%%%%%%%%%%%%%%%%%
\newcommand{\todo}[1]{\textcolor{red}{#1}}

%%% \chapter %%%%%%%%%%%%%%%%%%%%%%%%%%%%%%%%%%%%%%%%%%%%%%%%
\usepackage{titlesec}
\titleformat
{\chapter} % command
%[] % shape
{\bfseries\Huge} % format
{\thechapter} % label
{1.5ex} % sep
{} % before-code
%[] % after-code

% per settare noindent
\setlength{\parindent}{0pt}


% Crea un capitolo senza numerazione che pero` appare nell'indice %
\newcommand{\problemchapter}[1]{%
  \chapter*{#1}%
  \addcontentsline{toc}{chapter}{#1}%
\markboth{#1}{#1}
}

% Numerazione delle appendici secondo norma
\addto\appendix{
\renewcommand{\thesection}{\Alph{chapter}.\arabic{section}}
\renewcommand{\thesubsection}{\thesection.\arabic{subsection}}}

\setcounter{secnumdepth}{5} 	%per avere più livelli nei titoli
\setcounter{tocdepth}{5}		%per avere più livelli nell'indice


\titolo{Branch and Bound as PTAS}
\studente{Tiziano Zamaroni \vspace{1em}\\ Sergio Aquilini}
\relatore{Monaldo Mastrolilli}
\correlatore{-}
\committente{SUPSI}
\corso{C-P6051P - Seminari di Progetto}
\modulo{M-P6050P - Progetto di Semestre}
\anno{2025}

\begin{document}

\pagenumbering{alph}
\maketitle
\onehalfspacing
\frontmatter


\pagenumbering{roman}
\tableofcontents
\listoffigures					% Opzionale
\listoftables					% Opzionale

\newpage
\mainmatter
\pagenumbering{arabic}
\setcounter{page}{1}

\chapter{Introduzione}

\begin{itemize}
    \item Problemi NP-Hard
    \item Branch-and-Bound Algorithms as Polynomial-time Approximation Schemes
    \item Unrelated parallel machine scheduling problem
    \item Uniform/Identical parallel machine scheduling problem (identical caso particolare di uniform)
    \item Profiling come ulteriore ottimizzazione di identical machine scheduling
    \item Dimostrazione pratica della validità dell'approccio con profiling
\end{itemize}

\chapter{Conclusione}


\section{Risultati ottenuti}



\section{Possibili miglioramenti futuri}
Elenco di possibili sviluppi futuri:
\begin{itemize}
    \item abc
    \item def
\end{itemize}

\bibliographystyle{unsrt}
\bibliography{bibliografia}
\end{document}
